% 1
\frame[plain]{\titlepage}

% 2
\begin{frame}{Что мы уже умеем}
\small
На прошлой лекции мы научились доказывать, что при подходящем шаге градиентный спуск уменьшает функцию:
\[
 f(x_{k+1})<f(x_k).
\]

Но остаётся более фундаментальный вопрос:
\[
 \boxed{\text{К какому минимуму мы вообще можем прийти?}}
\]

\centering
\includegraphics[width=0.82\linewidth]{convex_vs_nonconvex.pdf}
\end{frame}

% 3
\begin{frame}{Проблема локального минимума}
\small
Для невыпуклой функции возможны разные стационарные точки:
\begin{itemize}
    \item локальные минимумы разного качества;
    \item локальные максимумы;
    \item седловые точки.
\end{itemize}

\begin{alertblock}{Главный вопрос}
Можно ли выделить класс функций, для которых \textbf{любой локальный минимум автоматически глобален}?
\end{alertblock}

Ответ этой лекции: да. Такой класс задаёт \textbf{выпуклость}.
\end{frame}

% 4
\begin{frame}{Почему выпуклость так полезна}
\small
\[
\boxed{
\text{выпуклость}
\Longrightarrow
\text{простая геометрия}
\Longrightarrow
\text{простые условия оптимальности}
\Longrightarrow
\text{глобальные гарантии}
}
\]

\vspace{0.7em}
Сегодня мы свяжем три уже знакомых объекта:
\[
\text{геометрию функции}
\qquad
\grad f(x)
\qquad
\grad^2 f(x).
\]

И увидим, как они позволяют распознавать глобальный минимум.
\end{frame}

% 5
\begin{frame}{Отрезок между двумя точками}
\small
Пусть $x,y\in\R^d$. Точки отрезка между ними имеют вид
\[
 z(t)=(1-t)x+ty,\qquad t\in[0,1].
\]

При $t=0$ получаем $x$, при $t=1$ --- $y$.

\begin{block}{Идея}
Выпуклое множество устроено так, что вместе с любыми двумя своими точками содержит и весь соединяющий их отрезок.
\end{block}
\end{frame}

% 6
\begin{frame}{Определение выпуклого множества}
\small
\begin{block}{Определение}
Множество $X\subseteq\R^d$ называется \textbf{выпуклым}, если для любых $x,y\in X$ и любого $t\in[0,1]$
\[
 \boxed{(1-t)x+ty\in X.}
\]
\end{block}

Это определение требует проверки \textbf{для любой пары точек} множества.
\end{frame}

% 7
\begin{frame}{Как выглядит выпуклость геометрически}
\centering
\includegraphics[width=0.96\linewidth]{convex_sets.pdf}

\vspace{0.25em}
\small
В невыпуклом множестве хотя бы один отрезок между двумя допустимыми точками выходит наружу.
\end{frame}

% 8
\begin{frame}{Самопроверка: какие множества выпуклы?}
\small
Определите, какие из множеств выпуклы:
\[
X_1=\{x:\norm{x}_2\le1\},
\]
\[
X_2=\{x\in\R^2:x_1+x_2\le1\},
\]
\[
X_3=\{x:\norm{x}_2=1\},
\]
\[
X_4=\{x\in\R^2:x_1x_2\ge1\}.
\]

\begin{alertblock}{Совет}
Чтобы показать невыпуклость, достаточно найти \textbf{одну} пару точек, для которой отрезок выходит из множества.
\end{alertblock}
\end{frame}

% 9
\begin{frame}{Выпуклая комбинация}
\small
Пусть $x_1,\dots,x_m\in\R^d$, а коэффициенты удовлетворяют
\[
 \alpha_i\ge0,
 \qquad
 \sum_{i=1}^m\alpha_i=1.
\]
Тогда
\[
 \sum_{i=1}^m\alpha_i x_i
\]
называется \textbf{выпуклой комбинацией} точек $x_i$.

\centering
\includegraphics[width=0.46\linewidth]{convex_combination.pdf}
\end{frame}

% 10
\begin{frame}{Выпуклое множество содержит выпуклые комбинации}
\small
\begin{block}{Утверждение}
Если $X$ выпукло и $x_1,\dots,x_m\in X$, то при
\[
\alpha_i\ge0,\qquad \sum_i\alpha_i=1
\]
имеем
\[
\boxed{\sum_i\alpha_i x_i\in X.}
\]
\end{block}

Для двух точек это ровно определение. Для трёх точек можно сначала объединить $x_1,x_2$, а затем полученную точку --- с $x_3$.
\end{frame}

% 11
\begin{frame}{Почему область определения тоже должна быть выпуклой}
\small
Дальше мы будем сравнивать значения функции в точках
\[
(1-t)x+ty.
\]
Если эта точка выходит из области определения, само сравнение теряет смысл.

\begin{block}{Поэтому}
Говоря о выпуклой функции $f:X\to\R$, обычно предполагают, что её область $X$ сама выпукла.
\end{block}
\end{frame}

% 12
\begin{frame}{Геометрическая идея выпуклой функции}
\centering
\includegraphics[width=0.92\linewidth]{chord_condition.pdf}

\vspace{0.15em}
\small
График выпуклой функции лежит \textbf{не выше} любой хорды, соединяющей две точки её графика.
\end{frame}

% 13
\begin{frame}{Определение выпуклой функции}
\small
\begin{block}{Определение}
Функция $f:X\to\R$ на выпуклом множестве $X$ называется \textbf{выпуклой}, если для любых $x,y\in X$ и $t\in[0,1]$
\[
\boxed{
 f((1-t)x+ty)
 \le
 (1-t)f(x)+tf(y).
}
\]
\end{block}
\end{frame}

% 14
\begin{frame}{Как читать определение буквально}
\small
Слева:
\[
 f\bigl(\underbrace{(1-t)x+ty}_{\text{выпуклая комбинация аргументов}}\bigr).
\]

Справа:
\[
 \underbrace{(1-t)f(x)+tf(y)}_{\text{та же комбинация значений функции}}.
\]

\begin{alertblock}{Смысл}
Сначала усреднить аргументы и затем вычислить функцию --- не хуже, чем сначала вычислить значения, а потом усреднить их.
\end{alertblock}
\end{frame}

% 15
\begin{frame}{Пример: $f(x)=x^2$ выпукла}
\footnotesize
Нужно проверить
\[
((1-t)x+ty)^2\le(1-t)x^2+ty^2.
\]
Рассмотрим разность правой и левой частей:
\[
\begin{aligned}
&(1-t)x^2+ty^2-((1-t)x+ty)^2\\
&=t(1-t)(x-y)^2\ge0.
\end{aligned}
\]

Следовательно, условие выпуклости выполняется.

\centering
\includegraphics[width=0.49\linewidth]{x2_chord.pdf}
\end{frame}

% 16
\begin{frame}{Невыпуклый пример}
\small
Функция
\[
 f(x)=x^3
\]
на всей $\R$ не является выпуклой.

Например, её кривизна меняет знак: слева от нуля
\[
 f''(x)=6x<0,
\]
а справа
\[
 f''(x)>0.
\]

Позже мы увидим, что это уже достаточная причина невыпуклости на всей прямой.
\end{frame}

% 17
\begin{frame}{Линейная функция — пограничный случай}
\small
Пусть
\[
 f(x)=a^\top x+b.
\]
Тогда
\[
\begin{aligned}
f((1-t)x+ty)
&=a^\top((1-t)x+ty)+b\\
&=(1-t)f(x)+tf(y).
\end{aligned}
\]

То есть в определении выпуклости всегда выполняется \textbf{равенство}.

\begin{block}{Следствие}
Линейная функция одновременно выпукла и вогнута.
\end{block}
\end{frame}

% 18
\begin{frame}{Неравенство Йенсена}
\small
Для выпуклой функции и выпуклой комбинации
\[
 \alpha_i\ge0,
 \qquad
 \sum_i\alpha_i=1
\]
справедливо
\[
\boxed{
 f\left(\sum_i\alpha_i x_i\right)
 \le
 \sum_i\alpha_i f(x_i).
}
\]

Это называется \textbf{неравенством Йенсена}.
\end{frame}

% 19
\begin{frame}{Йенсен — это обобщение определения}
\small
При $m=2$
\[
\alpha_1=1-t,
\qquad
\alpha_2=t,
\]
и неравенство Йенсена превращается в
\[
f((1-t)x+ty)
\le
(1-t)f(x)+tf(y).
\]

Для большего числа точек результат получается последовательным применением определения.
\end{frame}

% 20
\begin{frame}{Мини-задача на Йенсена}
\small
Функция $f(x)=e^x$ выпукла. Примените Йенсена к двум точкам с одинаковыми весами:
\[
\alpha_1=\alpha_2=\frac12.
\]

Получаем
\[
\boxed{
 e^{(x+y)/2}
 \le
 \frac{e^x+e^y}{2}.
}
\]

Это простейший пример того, как выпуклость автоматически порождает полезные неравенства.
\end{frame}

% 21
\begin{frame}{Но определение неудобно проверять напрямую}
\small
Чтобы проверить выпуклость по определению, формально нужно убедиться в неравенстве
\[
 f((1-t)x+ty)
 \le
 (1-t)f(x)+tf(y)
\]
для \textbf{всех} $x,y$ и \textbf{всех} $t\in[0,1]$.

\begin{alertblock}{Нужны локальные критерии}
Если функция дифференцируема, можно описать выпуклость через градиент. Если дважды дифференцируема --- через гессиан.
\end{alertblock}
\end{frame}
